\subsection{Non-Factorizable Projected Descriptions and Quantum Correlations}
  \label{subsec:non-injective-projection}

  A single admissible configuration of~$\chi$ may correspond to multiple
  distinct effective degrees of freedom.
  What appear as separate particles or subsystems are multiple projective
  manifestations of a single underlying ontological configuration.
  Effective degrees of freedom cannot be assigned independent states; any
  admissible description must be globally consistent with the
  underlying~$\chi$ configuration, leading to persistent correlations that
  do not rely on signal exchange or spacetime proximity.

  The deeper reason for this non-factorizability can be identified at the level of the
  spectral structure of the projection.
  A configuration $\omega \in \Omega_\chi$ may carry a global invariant $I(\omega)$---total
  charge, total angular momentum, winding number, or more generally any topological
  raccordement constraint---that belongs to a sector of the substrate which is not
  individually projectable within the effective description.
  Such invariants belong to spectral sectors of $L_\chi$ below the projection threshold
  $\lambda^*$ (Section~2.4) and therefore cannot appear as independent local observables
  in the effective description: they are not representable as a separable pair
  $I_A(\omega_A) \oplus I_B(\omega_B)$ projectable independently into $O_A \times O_B$,
  but exist once, at the level of $\chi$, and can only appear in the effective description
  as a property of the joint fiber $\Pi^{-1}(o_{AB})$.
  The relevant pre-image is therefore not $\Pi^{-1}(o_A) \times \Pi^{-1}(o_B)$ but the
  non-factorisable global class $\Pi^{-1}(o_{AB})$, whose elements are configurations
  constrained jointly by $I(\omega)$.
  This is the structural origin of entanglement: not a causal memory of past interaction,
  but the irreducible representation constraint imposed by a globally unique infra-projectable
  invariant.

  Spatial separation does not duplicate the invariant; it only changes the relational frame
  through which it is read.
  The two effective observables $o_A$ and $o_B$ are not two independent properties: they
  are two projections of the same substrate invariant, evaluated from distinct local relational
  contexts $\Pi_A$ and $\Pi_B$.
  Entanglement arises precisely in the regime where the effective description separates
  two subsystems that remain jointly constrained by an invariant which is not separately
  projectable in either subsystem alone.
  Measurement at $A$ selects a locally stable reprojection $\Pi_A(\omega) = o_A^*$; because
  the fiber $\Pi^{-1}(o_{AB})$ is non-factorizable, this selection simultaneously restricts the
  class of configurations admissible for $B$ to those compatible with $I(\omega)$ and $o_A^*$,
  without any signal propagating from $A$ to $B$.
  The apparent ``collapse'' at $B$ is a restriction of admissible projected descriptions
  imposed by the global uniqueness of the invariant, not a dynamical influence.

  \subsubsection*{Ontological Monism and the Shared Projection Hypothesis}
    \label{subsec:ontological-monism-shared-projection}

    A single connected excitation in~$\chi$ can be represented as several
    spatially separated effective excitations.
    The observed separation is a property of the projected metric
    representation, not a fundamental separation of the underlying entity.

    \paragraph{Shared Fiber Phase and Spin Correlations.}
      \label{subsec:shared-fiber-phase-spin}
      Spin correlations can be read as correlations of a \emph{shared}
      internal degree of freedom of the projection fiber.
      A measurement at one location selects a locally stable reprojection of
      that shared fiber degree of freedom; the correlated statistics at the
      distant location reflect the restricted set of reprojections still
      compatible with the same underlying configuration.

    \begin{figure}[t]
      \centering
      \begin{tikzpicture}[
        font=\small,
        node distance=10mm,
        box/.style={draw, rounded corners, align=center,
        inner sep=6pt},
        arrow/.style={-Latex, thick},
        note/.style={align=left, font=\footnotesize},
        micro/.style={draw, circle, inner sep=1.5pt}
      ]
    \node[micro] (c1) at (-2.4, 1.5) {};
    \node[micro] (c2) at (-1.2, 1.7) {};
    \node[micro] (c3) at ( 0.0, 1.55) {};
    \node[micro] (c4) at ( 1.2, 1.75) {};
    \node[micro] (c5) at ( 2.4, 1.6) {};
    \node[note, above=3mm of c3, align=center]
    {Multiple admissible $\chi$ configurations\\
    (same fiber under $\pi$)};
    \node[box, below=16mm of c3,
      minimum width=6.8cm] (chieff)
    {$\chi_{\mathrm{eff}}$\\
    \footnotesize single effective reality\\
    \footnotesize non-factorisable (entangled)};
    \node[box, below left=12mm and 14mm of chieff] (obsA)
    {Outcomes in context $A$};
    \node[box, below right=12mm and 14mm of chieff] (obsB)
    {Outcomes in context $B$};
    \draw[arrow] (c1) -- (chieff.north west);
    \draw[arrow] (c2) -- (chieff.north);
    \draw[arrow] (c3) -- (chieff.north);
    \draw[arrow] (c4) -- (chieff.north);
    \draw[arrow] (c5) -- (chieff.north east);
    \node[note, left=2mm of chieff, anchor=east]
    {\footnotesize non-injective\\
    \footnotesize projection $\pi$};
    \draw[arrow] (chieff)
    -- node[left=2mm, note] {$\mathcal{O}_A$} (obsA);
    \draw[arrow] (chieff)
    -- node[right=2mm, note] {$\mathcal{O}_B$} (obsB);
      \end{tikzpicture}
      \caption{Entangled (non-factorisable) regime.
      Multiple underlying $\chi$ configurations map to the same
        $\chi_{\mathrm{eff}}$.
        Different measurement contexts define distinct operational
        projections producing correlated outcomes without introducing
        multiple effective realities.}
      \label{fig:entanglement-two-projections}
      \end{figure}
